% GENERATED FILE. Edit canonical lessons, then run npm run generate:books.
% canonical-source-hash: fnv1a64:bd70257d83158186
% canonical-lessons: AR-C509-adwa, AR-C509-fayrus, AR-C509-jurthuma, AR-C509-tatim, AR-C509-waram

\chapter{Germs and jabs}
\label{ch:ar-germs-and-jabs}
\hlchaptermodality{509}

\begin{chapteropening}
\textbf{By the end of this chapter:} \emph{I can talk about an infection, a virus, a germ, a vaccination and a swelling.}

Complete the last lesson of chapter 509: I can talk about an infection, a virus, a germ, a vaccination and a swelling.
\end{chapteropening}

\section[\texorpdfstring{\ar{عدوى} (ʿadwā)}{عدوى (ʿadwā)}]{\ar{عدوى} (ʿadwā) — an infection, contagion}

\label{lesson:AR-C509-adwa}



\begin{quote}
Before the new one: say the Arabic for an inflammation, then the Arabic for to swell.
\end{quote}

\subsection*{You'll want to know: \ar{عدوى}}
\textbf{\ar{عدوى}} — \emph{ʿadwā} — \textquotedblleft{}an infection, contagion\textquotedblright{}.

\begin{quote}
\ar{انتقلت} \ar{العدوى} \ar{من} \ar{طفل} \ar{إلى} \ar{طفل}.
\end{quote}

\begin{quote}
\emph{intaqalat al-ʿadwā min ṭifl ilā ṭifl.}
\end{quote}

\begin{quote}
The infection passed from child to child.
\end{quote}

\begin{cousinweb}
\textbf{\ar{عدوى}} (\emph{ʿadwā}) is built on the native Arabic root \emph{ʿ-d-w}, \textquotedblleft{}to run, to pass beyond\textquotedblright{} (Lane). ʿAdwā is a thing that runs from one person to another.
\end{cousinweb}

\begin{grammarlens}[title={What you've built}]
One more word for describing what you lived and felt.
\end{grammarlens}

\subsection*{Your turn}
\begin{itemize}
\raggedright
  \item \emph{Say it:} \emph{ʿadwā}
  \item \emph{Say it:} \emph{ʿadwā}, once more
  \item \emph{Say it:} \emph{intaqalat al-ʿadwā min ṭifl ilā ṭifl.} — The infection passed from child to child.
  \item \emph{Recall:} say the Arabic for an inflammation, then the Arabic for to swell, then say \emph{ʿadwā} again
\end{itemize}

\begin{tcolorbox}[breakable,colback=teal!4,colframe=teal!35!black,title={Before you move on}]
What does \ar{عدوى} mean? (\textquotedblleft{}An infection, contagion\textquotedblright{}.) Say it once more.
\end{tcolorbox}
\section[\texorpdfstring{\ar{فيروس} (fayrūs)}{فيروس (fayrūs)}]{\ar{فيروس} (fayrūs) — a virus}

\label{lesson:AR-C509-fayrus}



\begin{quote}
Before the new one: say the Arabic for to swell, then the Arabic for an infection.
\end{quote}

\subsection*{You'll want to know: \ar{فيروس}}
\textbf{\ar{فيروس}} — \emph{fayrūs} — \textquotedblleft{}a virus\textquotedblright{}.

\begin{quote}
\ar{الزكام} \ar{سببه} \ar{فيروس}.
\end{quote}

\begin{quote}
\emph{az-zukām sababuhu fayrūs.}
\end{quote}

\begin{quote}
A cold is caused by a virus.
\end{quote}

\begin{cousinweb}
\textbf{\ar{فيروس}} (\emph{fayrūs}) is a modern loan, usually traced to Latin \emph{virus}, \textquotedblleft{}poison, slime\textquotedblright{}, which reached Arabic through European languages.
\end{cousinweb}

\begin{grammarlens}[title={What you've built}]
One more word for describing what you lived and felt.
\end{grammarlens}

\subsection*{Your turn}
\begin{itemize}
\raggedright
  \item \emph{Say it:} \emph{fayrūs}
  \item \emph{Say it:} \emph{fayrūs}, once more
  \item \emph{Say it:} \emph{az-zukām sababuhu fayrūs.} — A cold is caused by a virus.
  \item \emph{Recall:} say the Arabic for to swell, then the Arabic for an infection, then say \emph{fayrūs} again
\end{itemize}

\begin{tcolorbox}[breakable,colback=teal!4,colframe=teal!35!black,title={Before you move on}]
What does \ar{فيروس} mean? (\textquotedblleft{}A virus\textquotedblright{}.) Say it once more.
\end{tcolorbox}
\section[\texorpdfstring{\ar{جرثومة} (jurthūma)}{جرثومة (jurthūma)}]{\ar{جرثومة} (jurthūma) — a germ}

\label{lesson:AR-C509-jurthuma}



\begin{quote}
Before the new one: say the Arabic for an infection, then the Arabic for a virus.
\end{quote}

\subsection*{You'll want to know: \ar{جرثومة}}
\textbf{\ar{جرثومة}} — \emph{jurthūma} — \textquotedblleft{}a germ\textquotedblright{}.

\begin{quote}
\ar{يقتل} \ar{الصابون} \ar{الجراثيم}.
\end{quote}

\begin{quote}
\emph{yaqtulu aṣ-ṣābūn al-jarāthīm.}
\end{quote}

\begin{quote}
Soap kills germs.
\end{quote}

Plural: \ar{جراثيم}, \emph{jarāthīm}.

It is mostly met in the plural \ar{جراثيم}.

\begin{cousinweb}
\textbf{\ar{جرثومة}} (\emph{jurthūma}) is built on a native four-letter root, \emph{j-r-th-m}, where it named the root or origin of a thing; modern Arabic gave it the sense of a germ. The origin of the root itself is not settled.
\end{cousinweb}

\begin{grammarlens}[title={What you've built}]
One more word for describing what you lived and felt.
\end{grammarlens}

\subsection*{Your turn}
\begin{itemize}
\raggedright
  \item \emph{Say it:} \emph{jurthūma}
  \item \emph{Say it:} \emph{jurthūma}, once more
  \item \emph{Say it:} \emph{yaqtulu aṣ-ṣābūn al-jarāthīm.} — Soap kills germs.
  \item \emph{Recall:} say the Arabic for an infection, then the Arabic for a virus, then say \emph{jurthūma} again
\end{itemize}

\begin{tcolorbox}[breakable,colback=teal!4,colframe=teal!35!black,title={Before you move on}]
What does \ar{جرثومة} mean? (\textquotedblleft{}A germ\textquotedblright{}.) Say it once more.
\end{tcolorbox}
\section[\texorpdfstring{\ar{تطعيم} (taṭʿīm)}{تطعيم (taṭʿīm)}]{\ar{تطعيم} (taṭʿīm) — a vaccination}

\label{lesson:AR-C509-tatim}



\begin{quote}
Before the new one: say the Arabic for a virus, then the Arabic for a germ.
\end{quote}

\subsection*{You'll want to know: \ar{تطعيم}}
\textbf{\ar{تطعيم}} — \emph{taṭʿīm} — \textquotedblleft{}a vaccination\textquotedblright{}.

\begin{quote}
\ar{أخذ} \ar{الأطفال} \ar{التطعيم} \ar{في} \ar{المدرسة}.
\end{quote}

\begin{quote}
\emph{akhadha al-aṭfāl at-taṭʿīm fī al-madrasa.}
\end{quote}

\begin{quote}
The children had their vaccination at school.
\end{quote}

\begin{cousinweb}
\textbf{\ar{تطعيم}} (\emph{taṭʿīm}) is built on the native Arabic root \emph{ṭ-ʿ-m}, \textquotedblleft{}to eat, to taste\textquotedblright{} (Lane), the root of \emph{ṭaʿām}, food. Lane gives Form II \emph{ṭaʿʿama} as \textquotedblleft{}to graft\textquotedblright{} a branch; modern Arabic uses the same word for putting a vaccine into the body.
\end{cousinweb}

\begin{grammarlens}[title={What you've built}]
One more word for describing what you lived and felt.
\end{grammarlens}

\subsection*{Your turn}
\begin{itemize}
\raggedright
  \item \emph{Say it:} \emph{taṭʿīm}
  \item \emph{Say it:} \emph{taṭʿīm}, once more
  \item \emph{Say it:} \emph{akhadha al-aṭfāl at-taṭʿīm fī al-madrasa.} — The children had their vaccination at school.
  \item \emph{Recall:} say the Arabic for a virus, then the Arabic for a germ, then say \emph{taṭʿīm} again
\end{itemize}

\begin{tcolorbox}[breakable,colback=teal!4,colframe=teal!35!black,title={Before you move on}]
What does \ar{تطعيم} mean? (\textquotedblleft{}A vaccination\textquotedblright{}.) Say it once more.
\end{tcolorbox}
\section[\texorpdfstring{\ar{ورم} (waram)}{ورم (waram)}]{\ar{ورم} (waram) — a swelling; a tumour}

\label{lesson:AR-C509-waram}



\begin{quote}
Before the new one: say the Arabic for a germ, then the Arabic for a vaccination.
\end{quote}

\subsection*{You'll want to know: \ar{ورم}}
\textbf{\ar{ورم}} — \emph{waram} — \textquotedblleft{}a swelling; a tumour\textquotedblright{}.

\begin{quote}
\ar{ظهر} \ar{ورم} \ar{صغير} \ar{على} \ar{ركبتي}.
\end{quote}

\begin{quote}
\emph{ẓahara waram ṣaghīr ʿalā rukbatī.}
\end{quote}

\begin{quote}
A small swelling appeared on my knee.
\end{quote}

Plural: \ar{أورام}, \emph{awrām}.

\begin{cousinweb}
\textbf{\ar{ورم}} (\emph{waram}) is built on the native Arabic root \emph{w-r-m}, \textquotedblleft{}to swell\textquotedblright{} (Lane). Tawarrama, to swell up, is from the same root.
\end{cousinweb}

\begin{grammarlens}[title={What you've built}]
One more word for describing what you lived and felt.
\end{grammarlens}

\subsection*{Your turn}
\begin{itemize}
\raggedright
  \item \emph{Say it:} \emph{waram}
  \item \emph{Say it:} \emph{waram}, once more
  \item \emph{Say it:} \emph{ẓahara waram ṣaghīr ʿalā rukbatī.} — A small swelling appeared on my knee.
  \item \emph{Recall:} say the Arabic for a germ, then the Arabic for a vaccination, then say \emph{waram} again
\end{itemize}

\begin{tcolorbox}[breakable,colback=teal!4,colframe=teal!35!black,title={Before you move on}]
What does \ar{ورم} mean? (\textquotedblleft{}A swelling; a tumour\textquotedblright{}.) Say it once more.
\end{tcolorbox}
